%% ma: L10-B02-0002
%% lop: 10
%% chuong: 1
%% bai: 2
%% dang: 10.2.1
%% loai: TN4
%% muc: TH
%% dap_an: "C"
%% nguon: "Ảnh giáo viên gửi 10/10/2026 (ThS. Nguyễn Hoàng Việt, Chương 1 Mệnh đề), B-Đề 2, Phần 1 Câu 6"
%% kien_thuc: "Bài 2 (kí hiệu ∈, ⊂; tập con)"
%% trang_thai: cho-duyet
%% ly_do: "Dạng toán mới đề xuất 10.2.1: kí hiệu ∈, ⊂ giữa phần tử và tập hợp"
\begin{ex}
Cho biết $x$ là một phần tử của tập hợp $A$, xét các mệnh đề sau:

(I). $x\in A$. \quad (II). $\{x\}\in A$. \quad (III). $x\subset A$. \quad (IV). $\{x\}\subset A$.

Trong các mệnh đề sau, mệnh đề nào là đúng?
\choice{I và II}{I và III}{\True I và IV}{II và IV}
\loigiai{$x$ là phần tử của $A$ nên (I) $x\in A$ đúng; (IV) $\{x\}\subset A$ đúng (tập $\{x\}$ là tập con của $A$). (II) sai vì $\{x\}$ là tập hợp, không phải phần tử của $A$; (III) sai vì $\subset$ chỉ dùng giữa hai tập hợp. Chọn I và IV.}
\end{ex}
